\documentclass[11pt]{article}
\usepackage[margin=1.05in]{geometry}
\usepackage{amsmath, amssymb}
\usepackage{booktabs}
\usepackage{multirow}
\usepackage{xcolor}
\usepackage{hyperref}
\hypersetup{colorlinks=true, linkcolor=blue!50!black, urlcolor=blue!50!black}

\newcommand{\code}[1]{\texttt{#1}}
\DeclareMathOperator{\hn}{HN}
\DeclareMathOperator{\sd}{sd}

\title{\textbf{Time-block forecast scores: AR(2) versus i.i.d.\ temporal
pooling}\\[2pt]
\large Report on Experiments A--C (post \code{z.init} fix, assertion-guarded)}
\author{Score report}
\date{17 July 2026 \quad(rev.\ 18 July 2026: long-memory settings)}

\begin{document}
\maketitle

\begin{abstract}
Two clean experiments compare the AR(2) skew-t model against the
i.i.d.-in-time skew-t baseline under the time-block forecast protocol, by
CRPS and by the Brier score at the $q_{95}$ threshold, resolved by forecast
lead time. \textbf{Experiment A} (setting 4, five expanding-window blocks
$\times$ 5 datasets): the models tie once the reflected-ridge blocks are
excluded; the only visible AR(2) gain sits at lead 1 and is gone by lead 3,
exactly as the geometric memory of $\phi = (0.80, -0.35)$ predicts.
\textbf{Experiment B} (block 1 only, guarded fits, datasets 1--10): at the
same setting 4 the tie replicates; at the near-unit-root setting 5,
$\phi = (0.15, 0.80)$, the AR(2) advantage becomes real and large --- CRPS
$-31\%$ over leads 1--5 (Wilcoxon $p = 0.040$; sensitivity $p = 0.010$--$0.021$)
and a Brier improvement that is \emph{flat in lead time} (AR(2) better at 14
of 15 leads, mean $-0.008$), consistent with a memory horizon
($h^\ast \approx 109$ days) that dwarfs the 15-day window.
\textbf{Experiment C} (block 1, two long-memory settings, ARFIMA$(0,d,0)$ with
hyperbolic-decay dependence the AR(2) model cannot represent): the advantage
survives at the stronger $d = 0.45$ (Hurst $0.95$) --- CRPS $-13\%$, AR(2)
better at \emph{all} 15 Brier leads --- but is diluted by a short-window
identification failure that flags half the datasets, and vanishes at
$d = 0.35$. Across all four settings the advantage tracks the ratio of memory
horizon to training-window length, not the Hurst exponent, and remains
bounded by the score-ceiling
$\Delta\mathrm{CRPS}/\mathrm{CRPS} = 1 - \sqrt{1 - \mathrm{SNR}\,\kappa_z(h)}$.
\end{abstract}

\section{Data-generating processes and design}

All settings share one skeleton and differ \emph{only} in the temporal
dependence of the latent processes
(\code{code/analysis/time\_block\_forecast/simstudy/setup.R}):

\begin{center}
\small
\begin{tabular}{lp{0.72\textwidth}}
\toprule
component & specification \\
\midrule
marginal law & skew-$t$: $y_{st} = \mathbf{x}'\boldsymbol\beta + \lambda z_t
+ \varepsilon_{st}/\sqrt{\tau_t}$, $\tau \sim$ Gamma$(6, 16)$ mixing
($t$-tails, $\nu \approx 6$), $\lambda = 3$ (right skew), $\beta_0 = 10$ \\
skew latent & $z_t \mid \tau_t \sim \hn(0, \tau_t^{-1/2})$, global
($K = 1$ knot, so the temporal signal lives entirely in $\tau, z$) \\
spatial & Mat\'ern $\nu = 0.5$ (exponential), range $1$, spatial
proportion $\gamma = 0.9$; $144$ sites uniform on $[0,10]^2$ \\
temporal & Gaussian-copula \textbf{AR(2)} (settings 4--5) or
\textbf{ARFIMA$(0,d,0)$} (settings 6--7) on $\tau^\star, z^\star$ ($w^\star$),
per setting \\
record & $n_t = 200$ days; expanding-window blocks, each forecasting
$H = 15$ days \\
\bottomrule
\end{tabular}
\end{center}

\begin{center}
\small
\begin{tabular}{clll}
\toprule
id & label & temporal structure & persistence \\
\midrule
4 & strong          & AR(2) $\phi = (0.80, -0.35)$ & $\rho(F) = 0.592$,
$h^\ast = 6$ d (short memory) \\
5 & near-unit-root   & AR(2) $\phi = (0.15, 0.80)$ & $\rho(F) = 0.973$,
$h^\ast = 109$ d (short memory) \\
6 & LM $d = 0.35$    & ARFIMA$(0, 0.35, 0)$ & Hurst $0.85$, hyperbolic ACF
(long memory) \\
7 & LM $d = 0.45$    & ARFIMA$(0, 0.45, 0)$ & Hurst $0.95$, hyperbolic ACF
(long memory) \\
\bottomrule
\end{tabular}
\end{center}

\noindent Settings 4--5 are short-memory ($\sum_h |\rho(h)| < \infty$, ACF
$\sim \rho(F)^h$); settings 6--7 are genuinely long-memory
($\sum_h |\rho(h)| = \infty$, ACF $\sim h^{2d-1}$), which no finite-order AR
reproduces --- so the AR(2) model is \emph{misspecified by construction} there.
Experiments A and B use settings 4--5; Experiment C uses 6--7.

\paragraph{Protocol.} Fit on a contiguous prefix $y[, 1\!:\!T_o]$ (all sites
observed), forecast the window $(T_o, T_o + 15]$ with the seam-conditioned
recursion of \code{forecast\_block()}. Scores per lead $h$: CRPS
(\code{crps\_sample}, averaged over sites) and the Brier score for exceedance
of the validation block's $q_{95}$, identical thresholds for both methods,
paired by dataset. Both experiments post-date the \code{z.init} repair; all
Experiment-B fits passed the pre-registered consistency guard (truth
recovery $\wedge$ predictive spread; see the companion note
\code{tex/guard\_effect}), which replaced 13 of 40 attempt-0 chains.
Experiment A excludes the four blocks flagged by the same assertions
(reflected-ridge fits; \code{tex/lambda\_phiz\_ridge}).

\section{Experiment A: setting 4, five blocks $\times$ five datasets}

Expanding-window blocks with seams at $T_o = 50, 80, 110, 140, 170$;
$21$ clean (dataset, block) pairs.

\begin{table}[ht]
\centering
\caption{Experiment A, setting 4 (strong), clean 21 blocks. Mean score by
lead; gap $=$ AR(2) $-$ i.i.d.\ (negative favours AR(2)).}
\label{tab:expA}
\small
\begin{tabular}{crrrrrr}
\toprule
lead $h$ & \multicolumn{3}{c}{CRPS} & \multicolumn{3}{c}{Brier $q_{95}$} \\
\cmidrule(lr){2-4}\cmidrule(lr){5-7}
 & i.i.d. & AR(2) & gap & i.i.d. & AR(2) & gap \\
\midrule
1 & 1.822 & 1.514 & $-0.308$ & 0.0277 & 0.0190 & $-0.0087$ \\
2 & 2.120 & 2.048 & $-0.072$ & 0.0585 & 0.0552 & $-0.0034$ \\
3 & 2.273 & 2.288 & $+0.015$ & 0.0509 & 0.0493 & $-0.0016$ \\
4 & 2.933 & 2.940 & $+0.008$ & 0.0644 & 0.0644 & $+0.0000$ \\
5 & 3.330 & 3.318 & $-0.012$ & 0.0977 & 0.0989 & $+0.0012$ \\
6 & 2.785 & 2.941 & $+0.156$ & 0.0692 & 0.0801 & $+0.0109$ \\
7 & 2.285 & 2.261 & $-0.024$ & 0.0518 & 0.0523 & $+0.0006$ \\
8 & 2.611 & 2.576 & $-0.035$ & 0.0694 & 0.0701 & $+0.0007$ \\
9 & 1.968 & 1.991 & $+0.023$ & 0.0298 & 0.0301 & $+0.0003$ \\
10 & 2.150 & 2.062 & $-0.088$ & 0.0679 & 0.0603 & $-0.0076$ \\
11 & 2.055 & 2.085 & $+0.030$ & 0.0403 & 0.0417 & $+0.0014$ \\
12 & 2.372 & 2.520 & $+0.148$ & 0.0554 & 0.0608 & $+0.0053$ \\
13 & 3.087 & 3.232 & $+0.145$ & 0.0712 & 0.0767 & $+0.0055$ \\
14 & 2.350 & 2.267 & $-0.083$ & 0.0205 & 0.0231 & $+0.0026$ \\
15 & 2.432 & 2.298 & $-0.134$ & 0.0768 & 0.0718 & $-0.0050$ \\
\midrule
1--5 & 2.495 & 2.422 & $-0.074$ & 0.0599 & 0.0574 & $-0.0025$ \\
1--15 & 2.438 & 2.423 & $-0.015$ & 0.0568 & 0.0569 & $+0.0001$ \\
\bottomrule
\end{tabular}
\end{table}

\textbf{Reading.} The AR(2) gain is visible exactly where geometric memory
puts it: lead 1 ($-0.308$ CRPS, $-0.0087$ Brier, both the largest gaps in
the table) and residually lead 2; from lead 3 onward the gaps oscillate
around zero at the noise scale. Pooled over all leads the models tie
(paired two-sided tests over the 21 blocks: Brier $p = 0.97$, CRPS
$p = 0.78$). With spectral radius $0.592$, the seam state is worth
$\kappa_z(h) \approx 0.59^h$ of predictable signal: about a third at
$h = 1$ and effectively nothing by $h = 5$ --- the lead profile of
Table~\ref{tab:expA} is that decay, read through two proper scores.

\section{Experiment B: block 1, guarded, datasets 1--10}

Block 1 trains on $50$ days and forecasts days $51$--$65$. Every kept fit
passed the guard; setting-5 datasets $2$ and $8$ are excluded from the
primary analysis (flagged; sensitivity below).

\begin{table}[ht]
\centering
\caption{Experiment B, setting 4 (strong), $n = 10$ datasets.}
\label{tab:expB4}
\small
\begin{tabular}{crrrrrr}
\toprule
lead $h$ & \multicolumn{3}{c}{CRPS} & \multicolumn{3}{c}{Brier $q_{95}$} \\
\cmidrule(lr){2-4}\cmidrule(lr){5-7}
 & i.i.d. & AR(2) & gap & i.i.d. & AR(2) & gap \\
\midrule
1 & 2.550 & 2.275 & $-0.275$ & 0.0486 & 0.0552 & $+0.0066$ \\
2 & 3.169 & 3.141 & $-0.028$ & 0.0923 & 0.0922 & $-0.0001$ \\
3 & 1.749 & 1.766 & $+0.017$ & 0.0281 & 0.0236 & $-0.0045$ \\
4 & 1.861 & 1.841 & $-0.020$ & 0.0414 & 0.0398 & $-0.0016$ \\
5 & 1.996 & 2.011 & $+0.015$ & 0.0524 & 0.0540 & $+0.0015$ \\
6 & 2.244 & 2.259 & $+0.014$ & 0.0333 & 0.0351 & $+0.0019$ \\
7 & 2.628 & 2.660 & $+0.032$ & 0.0930 & 0.0938 & $+0.0008$ \\
8 & 1.775 & 1.827 & $+0.051$ & 0.0358 & 0.0390 & $+0.0031$ \\
9 & 1.649 & 1.685 & $+0.035$ & 0.0255 & 0.0262 & $+0.0007$ \\
10 & 2.172 & 2.143 & $-0.029$ & 0.0936 & 0.0915 & $-0.0020$ \\
11 & 1.676 & 1.664 & $-0.012$ & 0.0333 & 0.0347 & $+0.0013$ \\
12 & 2.132 & 2.096 & $-0.036$ & 0.0774 & 0.0762 & $-0.0011$ \\
13 & 2.589 & 2.543 & $-0.046$ & 0.1410 & 0.1368 & $-0.0043$ \\
14 & 1.805 & 1.808 & $+0.003$ & 0.0311 & 0.0324 & $+0.0013$ \\
15 & 2.067 & 2.058 & $-0.009$ & 0.0694 & 0.0682 & $-0.0013$ \\
\midrule
1--5 & 2.265 & 2.207 & $-0.058$ & 0.0526 & 0.0530 & $+0.0004$ \\
1--15 & 2.137 & 2.118 & $-0.019$ & 0.0597 & 0.0599 & $+0.0002$ \\
\bottomrule
\end{tabular}
\end{table}

\begin{table}[ht]
\centering
\caption{Experiment B, setting 5 (near-unit-root), clean $n = 8$ datasets.
AR(2) is better on CRPS at 12 of 15 leads and on Brier at \textbf{14 of 15}
leads.}
\label{tab:expB5}
\small
\begin{tabular}{crrrrrr}
\toprule
lead $h$ & \multicolumn{3}{c}{CRPS} & \multicolumn{3}{c}{Brier $q_{95}$} \\
\cmidrule(lr){2-4}\cmidrule(lr){5-7}
 & i.i.d. & AR(2) & gap & i.i.d. & AR(2) & gap \\
\midrule
1 & 2.120 & 1.531 & $-0.588$ & 0.0529 & 0.0437 & $-0.0092$ \\
2 & 3.366 & 1.879 & $-1.487$ & 0.0651 & 0.0554 & $-0.0097$ \\
3 & 2.685 & 2.119 & $-0.567$ & 0.0718 & 0.0719 & $+0.0001$ \\
4 & 2.788 & 1.552 & $-1.236$ & 0.0258 & 0.0172 & $-0.0087$ \\
5 & 2.516 & 2.172 & $-0.344$ & 0.0956 & 0.0876 & $-0.0080$ \\
6 & 2.717 & 1.979 & $-0.738$ & 0.0285 & 0.0184 & $-0.0101$ \\
7 & 2.551 & 2.181 & $-0.371$ & 0.0709 & 0.0563 & $-0.0146$ \\
8 & 2.845 & 2.121 & $-0.724$ & 0.0470 & 0.0384 & $-0.0086$ \\
9 & 2.600 & 2.508 & $-0.092$ & 0.0804 & 0.0706 & $-0.0098$ \\
10 & 3.786 & 2.736 & $-1.050$ & 0.1091 & 0.1036 & $-0.0055$ \\
11 & 1.713 & 1.849 & $+0.135$ & 0.0213 & 0.0120 & $-0.0093$ \\
12 & 3.225 & 2.402 & $-0.824$ & 0.0806 & 0.0746 & $-0.0061$ \\
13 & 1.732 & 1.869 & $+0.137$ & 0.0332 & 0.0260 & $-0.0072$ \\
14 & 2.739 & 2.347 & $-0.392$ & 0.0512 & 0.0443 & $-0.0069$ \\
15 & 1.889 & 2.005 & $+0.116$ & 0.0260 & 0.0180 & $-0.0081$ \\
\midrule
1--5 & 2.695 & 1.851 & $-0.844$ & 0.0623 & 0.0552 & $-0.0071$ \\
1--15 & 2.618 & 2.083 & $-0.535$ & 0.0573 & 0.0492 & $-0.0081$ \\
\bottomrule
\end{tabular}
\end{table}

\begin{table}[ht]
\centering
\caption{Experiment B paired one-sided tests ($H_1$: AR(2) better),
per-dataset units. ``Sens.'' keeps the flagged datasets ($n = 10$).}
\label{tab:tests}
\small
\begin{tabular}{llrrrr}
\toprule
setting & endpoint & gap & $t$ $p$ & Wilcoxon $p$ & sens.\ $t$/W $p$ \\
\midrule
4 & CRPS, leads 1--5        & $-0.058$  & $0.135$ & $0.154$ & --- \\
4 & Brier $q_{95}$, 1--5    & $+0.0004$ & $0.580$ & $0.541$ & --- \\
4 & CRPS, 1--15             & $-0.019$  & $0.216$ & $0.270$ & --- \\
4 & Brier $q_{95}$, 1--15   & $+0.0002$ & $0.546$ & $0.500$ & --- \\
\midrule
5 & CRPS, leads 1--5        & $-0.844$  & $0.066$ & $\mathbf{0.040}$ & $0.041$ / $\mathbf{0.021}$ \\
5 & Brier $q_{95}$, 1--5    & $-0.0071$ & $0.232$ & $0.417$ & $0.224$ / $0.305$ \\
5 & CRPS, 1--15             & $-0.535$  & $0.085$ & $\mathbf{0.021}$ & $0.043$ / $\mathbf{0.010}$ \\
5 & Brier $q_{95}$, 1--15   & $-0.0081$ & $0.119$ & $0.221$ & $0.124$ / $0.380$ \\
\bottomrule
\end{tabular}
\end{table}

\section{Experiment C: long-memory settings, block 1, guarded, datasets 1--10}

The two long-memory settings replace the AR(2) latent dynamics with
ARFIMA$(0, d, 0)$, whose autocorrelation decays hyperbolically
($\rho(h) \sim h^{2d-1}$) rather than geometrically. No finite-order AR
reproduces this, so the fitted AR(2) model is misspecified by construction;
the question is whether its geometric-memory pooling still beats i.i.d.\ when
the truth has heavier-than-geometric memory. Same block-1 protocol and guard.
The short training window interacts badly with long memory: the guard flags
$2$ of $10$ datasets at $d = 0.35$ and $5$ of $10$ at $d = 0.45$ (both methods,
$\lambda$ stable but under-identified in a $50$-day window), so the clean $n$
is smaller here --- itself a finding about fitting long memory on short records.

\begin{table}[ht]
\centering
\caption{Experiment C, setting 6 (LM $d = 0.35$, Hurst $0.85$), clean
$n = 8$.}
\label{tab:expC6}
\small
\begin{tabular}{crrrrrr}
\toprule
lead $h$ & \multicolumn{3}{c}{CRPS} & \multicolumn{3}{c}{Brier $q_{95}$} \\
\cmidrule(lr){2-4}\cmidrule(lr){5-7}
 & i.i.d. & AR(2) & gap & i.i.d. & AR(2) & gap \\
\midrule
1 & 1.788 & 1.575 & $-0.212$ & 0.0089 & 0.0214 & $+0.0125$ \\
2 & 2.227 & 2.196 & $-0.032$ & 0.0383 & 0.0447 & $+0.0063$ \\
3 & 1.521 & 1.742 & $+0.221$ & 0.0085 & 0.0128 & $+0.0043$ \\
4 & 2.412 & 2.540 & $+0.128$ & 0.0673 & 0.0691 & $+0.0018$ \\
5 & 2.386 & 2.172 & $-0.214$ & 0.0296 & 0.0272 & $-0.0024$ \\
6 & 2.758 & 2.696 & $-0.062$ & 0.0536 & 0.0548 & $+0.0012$ \\
7 & 1.881 & 1.840 & $-0.041$ & 0.0122 & 0.0144 & $+0.0022$ \\
8 & 2.942 & 2.780 & $-0.161$ & 0.0764 & 0.0740 & $-0.0024$ \\
9 & 2.546 & 2.540 & $-0.006$ & 0.0591 & 0.0617 & $+0.0025$ \\
10 & 1.732 & 1.770 & $+0.038$ & 0.0110 & 0.0126 & $+0.0016$ \\
11 & 2.938 & 2.804 & $-0.134$ & 0.0664 & 0.0629 & $-0.0034$ \\
12 & 2.989 & 3.030 & $+0.041$ & 0.0968 & 0.0981 & $+0.0014$ \\
13 & 3.232 & 3.227 & $-0.005$ & 0.0576 & 0.0589 & $+0.0013$ \\
14 & 3.344 & 3.239 & $-0.105$ & 0.1035 & 0.1008 & $-0.0027$ \\
15 & 2.737 & 2.751 & $+0.014$ & 0.0796 & 0.0829 & $+0.0033$ \\
\midrule
1--5 & 2.067 & 2.045 & $-0.022$ & 0.0305 & 0.0350 & $+0.0045$ \\
1--15 & 2.496 & 2.460 & $-0.035$ & 0.0513 & 0.0531 & $+0.0018$ \\
\bottomrule
\end{tabular}
\end{table}

\begin{table}[ht]
\centering
\caption{Experiment C, setting 7 (LM $d = 0.45$, Hurst $0.95$), clean
$n = 5$. AR(2) is better on CRPS at 13 of 15 leads and on Brier at
\textbf{all 15} leads.}
\label{tab:expC7}
\small
\begin{tabular}{crrrrrr}
\toprule
lead $h$ & \multicolumn{3}{c}{CRPS} & \multicolumn{3}{c}{Brier $q_{95}$} \\
\cmidrule(lr){2-4}\cmidrule(lr){5-7}
 & i.i.d. & AR(2) & gap & i.i.d. & AR(2) & gap \\
\midrule
1 & 1.750 & 1.633 & $-0.116$ & 0.0560 & 0.0394 & $-0.0166$ \\
2 & 2.257 & 1.960 & $-0.298$ & 0.0459 & 0.0229 & $-0.0230$ \\
3 & 2.223 & 1.659 & $-0.565$ & 0.0552 & 0.0302 & $-0.0250$ \\
4 & 2.013 & 1.738 & $-0.276$ & 0.0505 & 0.0319 & $-0.0186$ \\
5 & 3.119 & 2.856 & $-0.263$ & 0.0586 & 0.0394 & $-0.0191$ \\
6 & 2.640 & 2.213 & $-0.427$ & 0.0508 & 0.0490 & $-0.0018$ \\
7 & 2.287 & 1.946 & $-0.341$ & 0.0442 & 0.0262 & $-0.0179$ \\
8 & 4.400 & 3.527 & $-0.873$ & 0.1744 & 0.1494 & $-0.0250$ \\
9 & 2.692 & 2.479 & $-0.213$ & 0.0652 & 0.0440 & $-0.0212$ \\
10 & 2.970 & 2.957 & $-0.013$ & 0.1085 & 0.0914 & $-0.0171$ \\
11 & 3.137 & 3.178 & $+0.041$ & 0.1056 & 0.0876 & $-0.0180$ \\
12 & 2.731 & 2.705 & $-0.026$ & 0.0984 & 0.0832 & $-0.0152$ \\
13 & 2.467 & 2.447 & $-0.020$ & 0.0481 & 0.0344 & $-0.0137$ \\
14 & 2.487 & 2.451 & $-0.036$ & 0.0584 & 0.0406 & $-0.0178$ \\
15 & 3.249 & 3.343 & $+0.094$ & 0.1073 & 0.0938 & $-0.0135$ \\
\midrule
1--5 & 2.273 & 1.969 & $-0.304$ & 0.0532 & 0.0328 & $-0.0205$ \\
1--15 & 2.695 & 2.473 & $-0.222$ & 0.0751 & 0.0576 & $-0.0176$ \\
\bottomrule
\end{tabular}
\end{table}

\begin{table}[ht]
\centering
\caption{Experiment C paired one-sided tests ($H_1$: AR(2) better).}
\label{tab:testsC}
\small
\begin{tabular}{llrrrr}
\toprule
setting & endpoint & gap & $t$ $p$ & Wilcoxon $p$ & sens.\ $t$/W $p$ \\
\midrule
6 & CRPS, leads 1--5      & $-0.022$  & $0.336$ & $0.417$ & $0.296$ / $0.380$ \\
6 & Brier $q_{95}$, 1--5  & $+0.0045$ & $0.749$ & $0.472$ & $0.763$ / $0.541$ \\
\midrule
7 & CRPS, leads 1--5      & $-0.304$  & $0.067$ & $0.053$ & $0.076$ / $0.077$ \\
7 & Brier $q_{95}$, 1--5  & $-0.0205$ & $0.205$ & $0.500$ & $0.196$ / $0.658$ \\
7 & CRPS, 1--15           & $-0.222$  & $0.199$ & $0.394$ & $0.211$ / $0.419$ \\
7 & Brier $q_{95}$, 1--15 & $-0.0176$ & $0.185$ & $0.295$ & $0.203$ / $0.581$ \\
\bottomrule
\end{tabular}
\end{table}

\section{Reading the tables}

\paragraph{1. The advantage is a dose response --- in memory relative to the
window, not in the Hurst exponent.} Pooled CRPS gap (leads 1--5): $-0.058$
(setting 4, $\rho(F) = 0.59$), $-0.844$ (setting 5, $\rho(F) = 0.97$),
$-0.022$ (setting 6, $d = 0.35$), $-0.304$ (setting 7, $d = 0.45$). The order
is \emph{not} monotone in nominal long-memory strength: the short-memory
near-unit-root setting 5 beats both genuinely long-memory settings. What
orders the four is the memory horizon measured against the $50$-day training
window. Setting 5's $h^\ast \approx 109$ days makes the seam state highly
predictive across the whole $15$-day window; setting 7's hyperbolic ACF is
strong but the AR(2) model can only capture its geometric part, and the
$50$-day window under-identifies even that (see point 4). Setting 6's
short-lag ACF ($\rho(1) \approx 0.4$) sits in the same undetectable band as
setting 4. The ceiling $\Delta\mathrm{CRPS}/\mathrm{CRPS} = 1 - \sqrt{1 -
\mathrm{SNR}\,\kappa_z(h)}$ still bounds every case: the seam state is worth
only what the \emph{fitted} model's $\kappa_z(h)$ retains of it.

\paragraph{2. The lead profile matches the memory.} Setting 4
($h^\ast = 6$ d): the gain lives at leads 1--2 and dies by lead 3, in
\emph{both} experiments (Tables \ref{tab:expA}, \ref{tab:expB4}; lead-1 CRPS
gaps $-0.308$, $-0.275$ --- a clean cross-protocol replication). Settings 5
and 7 (long horizon): the gain does \emph{not} decay inside the window ---
setting 7's Brier gap is negative at \emph{all 15} leads (Table~\ref{tab:expC7})
and setting 5's at 14/15, because a long-memory latent state is still
informative $15$ days out.

\paragraph{3. CRPS detects what Brier cannot yet.} Where the advantage is
real (settings 5 and 7) the two scores agree in direction but CRPS reaches
significance first: setting 5 CRPS $W\,p = 0.021$--$0.040$, setting 7 CRPS
$W\,p = 0.053$ ($p = 0.077$ on all $10$), while both Brier endpoints stay
$p > 0.1$. The $q_{95}$ Brier scores one binary tail event per site-day and
carries an order of magnitude more sampling noise; at setting 7 its
dataset-level gap is also dominated by one dataset ($d\!=\!3$), so with clean
$n = 5$ the signed-rank test is mixed even though the lead-averaged gap is
negative at every lead. Powering Brier at these gaps needs $\sim\!30$--$50$
clean datasets (\code{run\_block1\_guarded.R --datasets=11:50}).

\paragraph{4. Long memory on a short window is an identification problem, not
just a modelling one.} Setting 7 flags half its datasets: a $50$-day window
can realise a nearly flat draw of a long-memory $z$ process, which
under-identifies $\lambda$ (stable but too small), and the AR(2)
misspecification compounds it. This is why setting 7, despite the strongest
nominal memory, does not beat setting 5: the advantage it \emph{could}
deliver is partly eaten by fitting difficulty on $50$ days. The direct test
of ``does long memory need a long window'' is therefore block 5 (train $170$
days), not more datasets --- the natural next experiment.

\paragraph{5. What the ties at settings 4 and 6 mean.} Neither is evidence
that AR(2) pooling is useless; each says the process memory is short relative
to the $15$-day window, so most scored leads lie beyond it and any pooled
endpoint dilutes its own signal. The informative settings are the ones whose
horizon exceeds the window (5 and 7), which is where the advantage appears.

\paragraph{Provenance.} Experiment A: \code{simstudy/results/4-*.RData}
(post-fix rerun), scored by \code{scores.R}; clean-block set from
\code{fit\_diag\_4.RData}. Experiments B and C:
\code{block1\_positive\_control/results/blk1-*.RData} (guarded driver),
scored by the same protocol; flagged datasets from \code{blk1\_diag} tables.
Guard counterfactual and the score-channel decomposition of the reflected
ridge: companion notes \code{tex/guard\_effect} and \code{tex/z\_init\_bug}.

\end{document}
